\chapter{Mathematical Demonstrations}
\label{app:mathematical_demonstrations}
This appendix provides supporting derivations for properties of \emph{FELICS} referenced in the main text. The demonstrations address the singleton special case, the treatment of independent concepts, and the preservation of the dataset-average causal effect.

\section*{Equality of Causal Effect between FELICS singleton-groups and Original Metric}
\begin{math_box}[label=ma:one_item_group]
    Let the correlation group consist of exactly one concept,
    \begin{equation}
        s=\{C_m\}.
    \end{equation}

    The payout coefficient introduced in Section~\ref{sec:shapley_correction} becomes

    \begin{equation}
        \gamma(x,s)
        =
        \frac{v(s)}
        {\sum_{C_i\in s}v(\{C_i\})}
        =
        \frac{v(\{C_m\})}
        {v(\{C_m\})}
        =1.
    \end{equation}

    Since the reference set used to compute the correction factor
    $\bar\gamma(x,s,T)$ is restricted to sets of identical treatment size,
    the only admissible reference size is likewise $|T|=1$. Consequently,

    \begin{equation}
        \bar\gamma(x,s,T)=1,
    \end{equation}

    and the normalized value function reduces to

    \begin{equation}
        v_{\mathrm{adj}}(T)
        =
        \frac{v(T)}
        {\bar\gamma(x,s,T)}
        =
        v(T).
    \end{equation}

    The adjusted causal effect is then obtained from the Shapley formulation

    \begin{equation}
        CE_{\mathrm{adj}}(x,C_m,s)
        =
        \sum_{T\subseteq s\setminus\{C_m\}}
        \frac{|T|!(|s|-|T|-1)!}{|s|!}
        \left[
            v_{\mathrm{adj}}(T\cup\{C_m\})
            -
            v_{\mathrm{adj}}(T)
            \right].
    \end{equation}

    Since $s=\{C_m\}$, the only admissible coalition is the empty set,
    $T=\emptyset$. The weighting coefficient therefore simplifies to

    \begin{equation}
        \frac{0!(1-0-1)!}{1!}=1,
    \end{equation}

    and using the convention $v_{\mathrm{adj}}(\emptyset)=0$,

    \begin{equation}
        CE_{\mathrm{adj}}(x,C_m,s)
        =
        v_{\mathrm{adj}}(\{C_m\}).
    \end{equation}

    Because no normalization is applied for singleton groups,

    \begin{equation}
        v_{\mathrm{adj}}(\{C_m\})
        =
        v(\{C_m\}),
    \end{equation}

    which is exactly the value function underlying the causal-effect definition of Matton et al.~\cite{mattonWalkTalkMeasuring2024}. Therefore,

    \begin{equation}
        CE_{\mathrm{adj}}(x,C_m,\{C_m\})
        =
        CE_{\mathrm{Matton}}(x,C_m)
    \end{equation}

    showing that \emph{FELICS} reduces exactly to the original CE every correlation group contains only a single concept. Because the definitions of the explanation implied effect and faithfulness (see Equations~\ref{eq:ee} and \ref{eq:faithfulness})  are unchanged, both metrics are equivalent in this scenario.
\end{math_box}

\section*{No Additional Causal Effect for Independent Concepts in FELICS}
\begin{math_box}[label=ma:independence_behavior]
    Assume a correlation group $s$ whose concepts are independent relative to the
    dataset. By definition of the normalization introduced in
    Section~\ref{sec:shapley_correction}, this corresponds to

    \begin{equation}
        \bar\gamma(x,s,s)=\gamma(x,s).
    \end{equation}

    The adjusted value of the grand coalition therefore becomes

    \begin{equation}
        v_{\mathrm{adj}}(s)
        =
        \frac{v(s)}
        {\gamma(x,s)}.
    \end{equation}

    Substituting the definition of the payout coefficient,

    \begin{equation}
        \gamma(x,s)
        =
        \frac{v(s)}
        {\sum_{C_i\in s}v(\{C_i\})},
    \end{equation}

    yields

    \begin{equation}
        v_{\mathrm{adj}}(s)
        =
        \sum_{C_i\in s}v(\{C_i\}).
    \end{equation}

    Hence, after normalization, the value assigned to the grand coalition equals the
    sum of the individual atomic intervention effects.

    Assume furthermore that every concept contributes independently to every
    coalition. Then the marginal contribution of a concept is identical in every
    coalition in which it appears,

    \begin{equation}
        v_{\mathrm{adj}}(T\cup\{C_m\})
        -
        v_{\mathrm{adj}}(T)
        =
        v(\{C_m\}),
        \qquad
        \forall\,T\subseteq s\setminus\{C_m\}.
    \end{equation}

    Since the Shapley weights form a partition of unity,

    \begin{equation}
        \sum_{T\subseteq s\setminus\{C_m\}}
        \frac{|T|!(|s|-|T|-1)!}{|s|!}
        =
        1,
    \end{equation}

    the adjusted causal effect reduces to

    \begin{equation}
        CE_{\mathrm{adj}}(x,C_m,s)
        =
        v(\{C_m\}).
    \end{equation}

    Finally, $v(\{C_m\})$ is precisely the KL-divergence-based intervention value
    used by Matton et al.~\cite{mattonWalkTalkMeasuring2024} to define the causal
    effect of an individual concept (see Equation~\ref{eq:ce}). Consequently,

    \begin{equation}
        CE_{\mathrm{adj}}(x,C_m,s)
        =
        CE_{\mathrm{Matton}}(x,C_m),
    \end{equation}

    showing that no additional causal effect is attributed to correlation groups
    whose concepts behave independently relative to the dataset.
\end{math_box}

\section*{CE Dataset-Average Preservation}
\begin{math_box}[label=ma:dataset_average_preservation]
    \textbf{Claim:} The dataset-level average of all Causal Effect values computed by FELICS is \emph{approximately} equal to that of Matton et al.'s original metric:
    \begin{equation}
        \frac{1}{N_C} \sum_{x \in X} \sum_{C_m \in C(x)} CE_{\text{adj}}(x, C_m, s) \approx \frac{1}{N_C} \sum_{x \in X} \sum_{C_m \in C(x)} v(\{C_m\}).
    \end{equation}
    Here, $N_C$ is the total number of concepts, $C(x)$ are the concepts in question $x$, and $s$ is the correlation group containing $C_m$.

    \smallskip
    \textbf{Intuition:} The payout-normalization step is designed to compensate for the \emph{soft ceiling effect} (Section~\ref{sec:kl_ceiling}), which systematically reduces joint effects $v(s)$ relative to the sum of singleton effects $\sum_{C_m \in s} v(\{C_m\})$. If this compensation is applied consistently across the dataset, the \emph{total causal influence} attributed to all concepts should remain approximately unchanged, preserving the overall scale of CE values.

    \smallskip
    \textbf{Simplifying Assumptions for Demonstration:} To illustrate why approximate equality is plausible, two idealized conditions will be relied on that do not hold universally but capture the \emph{typical} behavior of the reference-set construction:
    \begin{enumerate}
        \item \textbf{Representative Reference Sets:} The reference set $\mathcal{R}(x,s,s)$ (Equations~\ref{eq:R_1}--\ref{eq:R_3}) is treated as a \emph{representative sample} of all groups of size $|s|$ in the dataset. In practice, $\mathcal{R}(x,s,s)$ contains at most $5$ elements, drawn prioritized by proximity (same group $\rightarrow$ same question $\rightarrow$ same dataset). For large datasets, this sample is approximately representative of the population of size-$|s|$ groups.

        \item \textbf{Symmetry of Payout Coefficients:} The payout coefficients $\gamma(x',s')$ for groups of the same size are are assumed to be \emph{exchangeable}, i.e., their distribution is invariant to permutation. This allows us to treat $\bar\gamma(x,s,s)$ as an estimate of a global average $\bar\gamma_k$ for size-$k$ groups.
    \end{enumerate}
    The assumptions above are sufficient to justify the \emph{approximate} preservation expected to be observed in practice.

    \smallskip
    \textbf{Demonstration Under Assumptions:} For any group $s$, the Shapley \emph{efficiency axiom} gives
    \begin{equation}
        \label{eq:shapley_efficiency_approx}
        \sum_{C_m \in s} CE_{\text{adj}}(x, C_m, s) = v_{\text{adj}}(s) = \frac{v(s)}{\bar\gamma(x,s,s)}.
    \end{equation}
    Under Assumption~(1), $\bar\gamma(x,s,s) \approx \bar\gamma_k$, where $\bar\gamma_k$ is the average payout coefficient for all size-$k$ groups:
    \begin{equation}
        \bar\gamma_k := \frac{1}{|{\mathcal G}_k|} \sum_{s' \in {\mathcal G}_k} \frac{v(s')}{\sum_{C_m \in s'} v(\{C_m\})}.
    \end{equation}
    Here, ${\mathcal G}_k$ denotes the set of all size-$k$ groups (distinct from the graph $G$ in Section~\ref{sec:joint_counterfactual}). Substituting into Equation~\ref{eq:shapley_efficiency_approx}:
    \begin{align}
        \sum_{s \in {\mathcal G}_k} \sum_{C_m \in s} CE_{\text{adj}}(x, C_m, s) & = \sum_{s \in {\mathcal G}_k} \frac{v(s)}{\bar\gamma(x,s,s)} \approx \sum_{s \in {\mathcal G}_k} \frac{v(s)}{\bar\gamma_k}                                    \\
                                                                                & = \frac{1}{\bar\gamma_k} \sum_{s \in {\mathcal G}_k} v(s) = \frac{1}{\bar\gamma_k} \cdot \bar\gamma_k \sum_{s \in {\mathcal G}_k} \sum_{C_m \in s} v(\{C_m\}) \\
                                                                                & = \sum_{s \in {\mathcal G}_k} \sum_{C_m \in s} v(\{C_m\}).
    \end{align}
    The second equality follows from the definition of $\bar\gamma_k$: multiplying $\bar\gamma_k$ by the sum of all $v(s)$ for $s \in {\mathcal G}_k$ yields the sum of all $\sum_{C_m \in s} v(\{C_m\})$. Summing over all group sizes $k$:
    \begin{align}
        \sum_{x \in X} \sum_{s \in S(x)} \sum_{C_m \in s} CE_{\text{adj}}(x, C_m, s) & \approx \sum_{x \in X} \sum_{s \in S(x)} \sum_{C_m \in s} v(\{C_m\}) \\
                                                                                     & = \sum_{x \in X} \sum_{C_m \in C(x)} v(\{C_m\}).
    \end{align}
    Dividing by $N_C$ preserves the approximation, showing that the dataset-level averages are approximately equal. The error introduced by the simplifications diminishes as the dataset grows and the reference sets become more representative.
\end{math_box}
